\section{Limitations}
\label{sec:limitations}

\paragraph{Scale and realism.} The tasks are small gridworlds with hand-designed
channels. They are model organisms: useful because the true objective is known
exactly and the misspecification can be switched off, not because they resemble
the settings where reward hacking matters most. Claims transfer to frontier
systems only as hypotheses to test there.

\paragraph{One algorithm, one architecture.} Every agent is PPO with the same
small CNN. Whether hacking emerges, and when, depends on the optimizer and on
exploration; the dose--response thresholds we report are properties of this
agent, not of the reward functions alone.

\paragraph{Design choices were calibrated.} The network architecture and the
default channel magnitudes were chosen in single-seed pilot runs so that Oracle
PPO learns every task within the budget and the default channel is exploited
(Appendix~\ref{app:details}). That Proxy PPO underperforms Oracle PPO on $J_t$
is therefore a sanity check of the design, not a finding; the dose--response
experiment reports outcomes across magnitudes so that the calibration is
visible rather than hidden.

\paragraph{Hand-crafted monitors.} Our monitors are linear models and shallow
trees over at most \res{n-features-privileged} hand-written trajectory features.
Stronger monitors could do better, and monitors with access to an agent's
internals or reasoning are out of scope. Hack labels count uses of the channel,
not harm; we therefore also report harmful hacks (channel use without task
success), and the label thresholds are task-specific.

\paragraph{Training against a monitor.} Our monitor-penalty experiment uses one
penalty weight, one monitor channel and a fixed monitor. A monitor that is
retrained as the agent changes, or a penalty that is annealed, may behave
differently.
